\cleardoublepage
\phantomsection
\pdfbookmark{Compendio}{Compendio}
\begingroup
\let\clearpage\relax
\let\cleardoublepage\relax
\chapter*{Sommario}
\noindent
Nell'ambito dello sviluppo di sistemi \textit{embedded}, la verifica della corretta implementazione dei protocolli di comunicazione rappresenta un aspetto cruciale per garantirne l'affidabilità, in particolare nei settori \textit{automotive} e industriale.\\
La presente tesi descrive le attività svolte durante il mio \textit{stage} universitario, incentrato sullo sviluppo di un \textit{framework} di \textit{test} per dispositivi \textit{embedded} e sulla realizzazione di \textit{test} per la verifica di specifici protocolli di comunicazione.\\\\
Durante il primo periodo, di diverse settimane, mi sono dedicato allo sviluppo, in linguaggio Python, di un \textit{framework} per l'esecuzione automatizzata di test, che ha comportato sia il \textit{refactoring} di \textit{test} preesistenti sia la progettazione e la scrittura di nuovi casi di test.
La maggior parte dei \textit{test} realizzati era volta a verificare la comunicazione su canale CAN (\textit{Controller Area Network}) e le sessioni diagnostiche UDS (\textit{Unified Diagnostic Services}) su CAN, protocollo largamente diffuso nella diagnostica \textit{automotive}.\\\\
Successivamente, ho esteso il progetto includendo la comunicazione Modbus su interfaccia UART, per la quale ho sviluppato \textit{test} volti a verificare quali \textit{function code} risultassero eseguibili dal dispositivo e su quali indirizzi.
Per rendere possibile questa verifica è stato necessario, in uno stadio preliminare, programmare in linguaggio C una scheda Nucleo-H533RE, affinché fosse in grado di rispondere alle richieste Modbus e di implementare correttamente i \textit{function code} previsti dal protocollo.
\newpage
\noindent
Nella tesi approfondisco inoltre le principali tecnologie e gli strumenti impiegati, discuto l'utilità dei \textit{test} sviluppati e illustro le motivazioni alla base delle soluzioni progettuali adottate.
Una sezione conclusiva è infine dedicata alle competenze acquisite durante l'esperienza, tra cui : competenze di programmazione in Python e in C, applicazione in ambito reale di concetti di progettazione del \textit{software} e acquisizione di conoscenze sui protocolli CAN, UDS e Modbus.

\endgroup
\vfill
